\section{Sharp countermodels and stop rules}
\label{sec:countermodels}

The following examples prevent five tempting but false extensions.

\begin{example}[An analysis leaf is not free]\label{ex:row-leaf}
Let
\[
 C=\begin{pmatrix}1\\-1\end{pmatrix},
 \qquad
 b=\begin{pmatrix}1\\1\end{pmatrix}.
\]
Both analysis rows are leaves, but
\[
 \min_{\theta\in\C}
 \bigl(|1+\theta|+|1-\theta|\bigr)=2.
\]
For \(b=(1,-1)^\mathsf T\), the value is instead zero.  The same tree
topology therefore implies neither a positive cost nor a saving for a
specific datum.
\end{example}

\begin{example}[A complex star is a geometric median]\label{ex:median}
For one variable with points \(0,1,\mathrm i\), the objective is
\[
 |\theta|+|\theta-1|+|\theta-\mathrm i|.
\]
The coordinatewise median is \(0\) and has value \(2\).  The planar
Fermat point
\[
 \theta_*=\frac{3-\sqrt3}{6}(1+\mathrm i)
\]
has value \(\sqrt{2+\sqrt3}<2\).  Complex star components are weighted
geometric-median problems; only collinear data reduce to an ordinary
weighted median.  With two points \(p_1,p_2\) and weights \(w_1,w_2\),
the exact value is
\[
 \min(w_1,w_2)|p_1-p_2|.
\]
\end{example}

\begin{example}[Cycle support does not fix rank]\label{ex:cycle-rank}
The matrices
\[
 C_+=\begin{pmatrix}1&1\\1&1\end{pmatrix},
 \qquad
 C_-=\begin{pmatrix}1&1\\1&-1\end{pmatrix}
\]
have the same \(K_{2,2}\) support, but ranks one and two.  For
\(b=(1,-1)^\mathsf T\),
\[
 \Gamma(b,C_+)=2,
 \qquad
 \Gamma(b,C_-)=0.
\]
The forest hypothesis in \cref{thm:forest-rank} is therefore necessary.
\end{example}

\begin{example}[Nonzero real-linear leaves need not be surjective]
Let a single real-linear block be \(C(z)=\Real z\), regarded as a map
\(\C\to\C\), and take \(b=\mathrm i\).  Then
\[
 \min_{z\in\C}|\mathrm i+\Real z|=1.
\]
The one-edge graph is a tree, but the leaf cannot be erased.  General
real-linear blocks require a surjective leaf map; the matching and
phase-transport formulas require further invertibility and conformality.
\end{example}

\begin{example}[Real and complex parameter fields differ]
\[
 \min_{t\in\R}|\mathrm i+t|=1,
 \qquad
 \min_{t\in\C}|\mathrm i+t|=0.
\]
For real parameters the dual stationarity condition is a real
projection equation, not \(C^*z=0\) as a complex equation.  No theorem
may switch parameter fields silently.
\end{example}

\begin{remark}[Topology is not a lower bound by itself]
A Hall circuit supplies a one-dimensional space of possible charges.
Its score can still vanish because
\(\langle g^S,b_S\rangle=0\).  Likewise, many motifs can overlap so
strongly that their fractional packing has small total capacity.
Arithmetic usefulness needs both carrier prevalence and noncancelling
literal scores.
\end{remark}

